\subsection{Transforming and Filtering Iteration Streams}

Traversal does not always mean merely receiving values exactly as they appear in the original source. Sometimes each value must be transformed before use. Sometimes only some values should continue through the pipeline. Python provides \texttt{map()} and \texttt{filter()} for these two common cases.

The important point is that both functions are lazy in Python 3. They do not immediately build a list of results. They return iterator objects that produce values only when the program asks for them.

\begin{quote}
\texttt{map()} transforms values lazily. \texttt{filter()} keeps or rejects values lazily.
\end{quote}

This makes them part of the same traversal architecture as \texttt{enumerate()} and \texttt{zip()}. They are not containers. They are traversal objects.

\subsubsection{\texttt{map()} and \texttt{filter()} as Lazy Iterator-Producing Transformations}

The \texttt{map()} function applies a callable object to each value produced by an iterable. Its basic shape is:

\begin{verbatim}
map(callable_object, iterable)
\end{verbatim}

A simple example uses the built-in \texttt{int} constructor to transform strings into integers:

%<<<PROGRAM:programs/tra03>>>
The \texttt{map} object is an iterator. It has \texttt{\_\_iter\_\_()} and \texttt{\_\_next\_\_()}. Therefore it is consumed as it advances.

Manual traversal shows the transformation:

\begin{verbatim}
texts = ["10", "20", "42"]

nums = map(int, texts)

print(f"next(nums): {next(nums)}")
print(f"next(nums): {next(nums)}")
print(f"next(nums): {next(nums)}")
\end{verbatim}

Possible output:

\begin{verbatim}
next(nums): 10
next(nums): 20
next(nums): 42
\end{verbatim}

Each original string is passed to \texttt{int}. The produced value is an integer object.

%<<<PROGRAM:programs/tra02>>>
The \texttt{\_\_add\_\_()} method supports integer addition. The \texttt{\_\_bool\_\_()} method supports truth-value testing.

A \texttt{map} object can be used directly in a \texttt{for} loop:

\begin{verbatim}
texts = ["10", "20", "42"]

for num in map(int, texts):
    print(f"{num:>3} | {num * 2:>3}")
\end{verbatim}

Possible output:

\begin{verbatim}
 10 |  20
 20 |  40
 42 |  84
\end{verbatim}

The loop receives one transformed value at a time. No list of transformed values is required unless the programmer explicitly asks for one.

Materialization creates stored results:

\begin{verbatim}
texts = ["10", "20", "42"]

nums = list(map(int, texts))

print(f"nums: {nums}")
print(f"type(nums): {type(nums)}")
\end{verbatim}

Possible output:

\begin{verbatim}
nums: [10, 20, 42]
type(nums): <class 'list'>
\end{verbatim}

The call to \texttt{list()} is the materialization boundary. Before that point, the result is a lazy \texttt{map} iterator. After that point, the result is a list containing the transformed values.

A \texttt{map} object is consumed after traversal:

%<<<PROGRAM:programs/tra01>>>
The first \texttt{list(nums)} consumes the iterator. The second call receives no remaining values.

To repeat the transformation, create a new \texttt{map} object:

\begin{verbatim}
texts = ["10", "20", "42"]

first = list(map(int, texts))
second = list(map(int, texts))

print(f"first: {first}")
print(f"second: {second}")
\end{verbatim}

Possible output:

\begin{verbatim}
first: [10, 20, 42]
second: [10, 20, 42]
\end{verbatim}

The source list is reusable. The \texttt{map} iterator is not reusable after exhaustion.

The transformation callable does not have to be \texttt{int}. It can be another callable object. A simple named function makes the lazy behavior visible:

%<<<PROGRAM:programs/tra06>>>


The function is not called when the \texttt{map} object is created. It is called when the iterator is advanced.

The function object itself is a runtime object:

\begin{verbatim}
def loud_text(text):
    return text.upper()

print(f"type(loud_text): {type(loud_text)}")
print(f"dir(loud_text) sample: {dir(loud_text)[:12]}")
print(f"loud_text.__name__: {loud_text.__name__}")
\end{verbatim}

Possible output:

\begin{verbatim}
type(loud_text): <class 'function'>
dir(loud_text) sample: ['__annotations__', '__builtins__', '__call__',
'__class__', '__closure__', '__code__', '__defaults__',
'__delattr__', '__dict__', '__dir__', '__doc__', '__eq__']
loud_text.__name__: loud_text
\end{verbatim}

The important method here is \texttt{\_\_call\_\_()}. A callable object can be called with parentheses. \texttt{map()} stores that callable and applies it to incoming values one at a time.

The \texttt{filter()} function has a related but different role. It keeps only the values for which a predicate succeeds.

Its basic shape is:

\begin{verbatim}
filter(predicate, iterable)
\end{verbatim}

The predicate is called for each value. If the predicate result is true, the original value is yielded. If the predicate result is false, the value is skipped.

%<<<PROGRAM:programs/tra05>>>


The \texttt{filter} object is also an iterator. It is lazy and consumable.

Using it in a loop:

\begin{verbatim}
def is_even(num):
    return num % 2 == 0

nums = [10, 15, 20, 42, 99]

for num in filter(is_even, nums):
    print(f"accepted: {num}")
\end{verbatim}

Possible output:

\begin{verbatim}
accepted: 10
accepted: 20
accepted: 42
\end{verbatim}

The values \texttt{15} and \texttt{99} were tested but not yielded, because the predicate returned false for them.

A manual version makes the logic explicit:

%<<<PROGRAM:programs/tra04>>>


The \texttt{filter()} version packages this pattern into a lazy iterator.

The predicate result is truth-tested, not required to be exactly \texttt{True} or \texttt{False}. This means \texttt{filter(None, iterable)} keeps truthy values and rejects falsy values.

\begin{verbatim}
values = ["spam", "", "Ni!", 0, 42, [], [1, 2]]

kept = filter(None, values)

print(f"type(kept): {type(kept)}")
print(f"list(kept): {list(kept)}")
\end{verbatim}

Possible output:

\begin{verbatim}
type(kept): <class 'filter'>
list(kept): ['spam', 'Ni!', 42, [1, 2]]
\end{verbatim}

The empty string, zero, and empty list were rejected because they are false in truth-value testing.

This connects \texttt{filter()} directly to the earlier truth-value rules:

\begin{verbatim}
values = ["spam", "", 0, 42, []]

for value in values:
    print(f"{repr(value):>8} -> bool: {bool(value)}")
\end{verbatim}

Possible output:

\begin{verbatim}
  'spam' -> bool: True
      '' -> bool: False
       0 -> bool: False
      42 -> bool: True
      [] -> bool: False
\end{verbatim}

The \texttt{filter(None, values)} form keeps only those values whose \texttt{bool()} result is true.

Like \texttt{map}, a \texttt{filter} object is consumed after traversal:

%<<<PROGRAM:programs/tra03>>>


The first materialization consumes the iterator.

The two functions can be combined into a traversal pipeline:

\begin{verbatim}
texts = ["10", "15", "20", "42"]

nums = map(int, texts)
even_nums = filter(lambda num: num % 2 == 0, nums)

print(list(even_nums))
\end{verbatim}

Possible output:

\begin{verbatim}
[10, 20, 42]
\end{verbatim}

This works, but the \texttt{lambda} expression introduces another compact function syntax. For early teaching code, a named function is often clearer:

%<<<PROGRAM:programs/tra02>>>


The pipeline is:

\begin{verbatim}
texts
    -> map int over each text
    -> filter only even integers
    -> materialize as list
\end{verbatim}

The important detail is that \texttt{nums} is itself a \texttt{map} iterator. The \texttt{filter} object consumes it as needed. Values pass through the chain one at a time.

A final example makes the stream-like behavior visible:

%<<<PROGRAM:programs/tra01>>>


The pipeline stops as soon as \texttt{next(large\_nums)} receives one accepted value. The value \texttt{"100"} has not yet been mapped or filtered, because no one has asked for another accepted value.

The practical rule is:

\begin{quote}
Use \texttt{map()} when every traversed value should be transformed. Use \texttt{filter()} when some traversed values should be allowed through and others rejected. Remember that both return lazy iterators, not stored result lists.
\end{quote}